\section{SmallWood Programming Model}
\label{sec:smallwood-model}
\label{sec:smallwood}

This section gives the SmallWood formulation of the ARC-SPRUCE NIZKs.  The prover commits to packed
witness rows, represented as low-degree row polynomials over a packing set $\Omega$ and masked on a
pairwise-disjoint hiding set $\Omega'$.  The verifier then checks PACS-style constraints on those committed
rows.  We record here the semantic statement layers and the exact coefficient-domain consequences needed by
the security arguments; concrete row inventories, gadget-level realizations, and proof-stack engineering are
deferred to Appendix~\ref{app:smallwood}.

\subsection{Proof interface}
\label{subsec:smallwood-interface}
Fix a packing set $\Omega=\{\omega_1,\ldots,\omega_s\}\subset\Fq$ and a disjoint hiding set
$\Omega'=\{\omega'_1,\ldots,\omega'_\ell\}\subset\Fq\setminus\Omega$.  A compiled witness is a matrix
\[
W\in\Fq^{n\times s}.
\]
For each row $i$, the prover interpolates a polynomial $P_i\in\Fq[X]$ of degree $<s+\ell$ such that
$P_i(\omega_j)=W_{i,j}$ for all $\omega_j\in\Omega$, while the values on $\Omega'$ are random masks.  The
committed witness is therefore a family of row polynomials
\[
P=(P_1,\ldots,P_n).
\]
SmallWood proves families of constraints over these rows.  Parallel constraints are slotwise identities
\[
  f_a\bigl(P_1(\omega),\ldots,P_n(\omega),\Theta_{a,1}(\omega),\ldots,\Theta_{a,t_a}(\omega)\bigr)=0
  \qquad\text{for every }\omega\in\Omega,
\]
where $\Theta_{a,j}$ are public helper rows.  Aggregated constraints are domain-summed identities
\[
  \sum_{\omega\in\Omega}
  g_b\bigl(P_1(\omega),\ldots,P_n(\omega),\Psi_{b,1}(\omega),\ldots,\Psi_{b,u_b}(\omega)\bigr)=0.
\]
In ARC-SPRUCE, the Ajtai commitment equation, centering equations, and BB-tran target identity are affine
coefficient-wise constraints; the inverse-witness equation
\((B_3-\mathbf 1_n r_1)\Had Z=1\) is bilinear; and the PRF checks contain the nonlinear Poseidon2-style
S-box constraints.  Range, carry, packing-consistency, and norm gadgets may additionally use aggregated
checks.

\subsection{Issuance as a SmallWood statement}
\label{subsec:smallwood-issuance}
\label{subsec:presign-compiled}
The public issuance statement is
\[
  x^{\mathrm{SW}}_{\mathrm{Issue}}=(\ComPublicParam,\vecc,r_0^I,r_1^I,T,B),
\]
and the hidden witness is
\[
  w^{\mathrm{SW}}_{\mathrm{Issue}}=(\vecmu,r_0^H,r_1^H,\bar r,r_0,r_1,Z).
\]
The committed witness rows encode the message lanes, the two holder-side randomness blocks, the commitment-only
randomness block, the long centered target-randomizer $r_0\in\Rq^{\ell_{r_0}}$, the inverse-side randomness $r_1$, the inverse witness $Z$, and the auxiliary carry rows used
by the centering gadget.  Semantically, the proof certifies the following claims.

\paragraph{Commitment opening.}
The public commitment opens to the same hidden message and holder-side randomness used later by the target:
\[
  \vecc=\mACom[\vecmu\Vert r_0^H\Vert r_1^H\Vert \bar r]^\top.
\]
Thus the message and holder randomness are fixed before the issuer sends its challenge randomness.

\paragraph{Centering constraints.}
The final BB-tran randomness is obtained from the committed holder randomness and the public issuer randomness:
\[
  r_0=\centerfunc(r_0^H+r_0^I),
  \qquad
  r_1=\centerfunc(r_1^H+r_1^I).
\]
In the compiled constraint system this is realized coefficient-wise through carry variables
$\kappa_0,\kappa_1\in\{-1,0,1\}$ satisfying
\[
  r_b^H+r_b^I-r_b=\kappa_b\cdot(2\BoundMsg+1)
  \qquad\text{for }b\in\{0,1\}.
\]
These carry rows are range-constrained and packed together with the source rows.

\paragraph{Inverse-witness equation.}
The proof certifies that $r_1$ is admissible and that $Z$ is its inverse witness:
\[
  (B_3-\mathbf 1_n r_1)\Had Z=1.
\]
This is the algebraic form used by the issuance proof; admissibility is thus enforced internally by the witness
rather than merely checked externally.

\paragraph{Target identity.}
The public target is the BB-tran target corresponding to the same hidden message and centered randomness:
\[
  T=B_0+B_1\vecmu+B_2r_0+Z.
\]
Hence the pre-sign proof binds one coherent hidden witness simultaneously to the Ajtai commitment, the
centered-randomness derivation, and the public target signed by the issuer.

\subsection{Showing as a SmallWood statement}
\label{subsec:smallwood-showing}
\label{subsec:showing-compiled}
The public showing statement is
\[
  x^{\mathrm{SW}}_{\mathrm{Show}}=(\nonce,\prftag,\mA,B,\BoundMsg,\BoundSig),
\]
with verifier state handled outside the proof.  The hidden witness is
\[
  w^{\mathrm{SW}}_{\mathrm{Show}}=(\vecu,\vecm,\veck,r_0,r_1,Z).
\]
The committed rows encode the short signature witness $\vecu$, the hidden message/key split $(\vecm,\veck)$,
the BB-tran randomness $(r_0,r_1)$ with $r_0\in\Rq^{\ell_{r_0}}$, the inverse witness $Z$, and the PRF companion rows.

Semantically, the proof certifies the following coefficient-domain statement:
\begin{align}
  (B_3-\mathbf 1_n r_1)\Had Z&=1, \label{eq:sw-show-inv}\\
  \mA\vecu&=B_0+B_1(\vecm\Vert\veck)+B_2r_0+Z, \label{eq:sw-show-sig}\\
  \prftag&=\PRF(\veck,\nonce), \label{eq:sw-show-prf}
\end{align}
together with the coefficient bounds on $\vecm,\veck,r_0,r_1$ and the shortness bound on $\vecu$.
Equation~\eqref{eq:sw-show-sig} is the direct witness-form BB-tran signature relation.  Eliminating $Z$ yields
\begin{equation}
  (B_3-\mathbf 1_n r_1)\Had\bigl(\mA\vecu-B_0-B_1(\vecm\Vert\veck)-B_2r_0\bigr)=1,
  \label{eq:sw-show-eliminated}
\end{equation}
which is the coefficient-domain relation used later in the security arguments.

\subsection{Replay / transform implementation}
\label{subsec:replay-implementation}
Some implementations check the showing relation in a transform or replay basis.  Let
\[
  \mathcal F_{\mathrm{rep}}:\Rq\to\mathcal T
\]
be the public replay map, with block projections $\mathcal F_{\mathrm{rep},b}$.  The prover then carries replay rows
for the images of
\[
  \mA\vecu,
  \qquad
  \vecm\Vert\veck,
  \qquad
  r_0,
  \qquad
  r_1.
\]
The compiled relation enforces exact source-to-replay bridges
\[
  \widehat T_b=\mathcal F_{\mathrm{rep},b}(\mA\vecu),
  \qquad
  \widehat\mu_b=\mathcal F_{\mathrm{rep},b}(\vecm\Vert\veck),
\]
\[
  \widehat r_{0,b}=\mathcal F_{\mathrm{rep},b}(r_0),
  \qquad
  \widehat r_{1,b}=\mathcal F_{\mathrm{rep},b}(r_1),
\]
and checks the replay residual
\begin{equation}
  (B_{3,b}^{\Theta}-\widehat r_{1,b})\Had
  \bigl(\widehat T_b-B_{0,b}^{\Theta}-B_{1,b}^{\Theta}\Had \widehat\mu_b-B_{2,b}^{\Theta}\Had\widehat r_{0,b}\bigr)
  =\mathcal F_{\mathrm{rep},b}(1).
  \label{eq:sw-replay-residual}
\end{equation}
Here $B_{i,b}^{\Theta}$ denotes the public replay image of the relevant BB-tran constant or coefficient block.

\begin{assumption}[Exact replay map]\label{ass:replay-transform}
The public replay map $\mathcal F_{\mathrm{rep}}$ used by the showing implementation is additive,
multiplicative on all products used in the compiled residual, and injective on the coefficient-domain values
that appear in the showing relation.  If a concrete replay representation is not injective, the compiled relation
must include explicit linking constraints sufficient to recover the same coefficient-domain identities.
\end{assumption}

\begin{theorem}[Replay-space certification implies the coefficient-domain BB-tran identity]
\label{thm:replay-to-coeff}
Assume Assumption~\ref{ass:replay-transform}.  Assume further that the compiled showing relation enforces
the exact bridges above and the replay residual \eqref{eq:sw-replay-residual}.  Then the extracted
coefficient-domain witness satisfies
\[
  (B_3-\mathbf 1_n r_1)\Had\bigl(\mA\vecu-B_0-B_1(\vecm\Vert\veck)-B_2r_0\bigr)=1
\]
in $\Rq$.
\end{theorem}

\begin{proof}
Substitute the exact bridge equalities into the replay residual.  Additivity and multiplicativity of
$\mathcal F_{\mathrm{rep}}$ transform the replay residual into the replay image of the coefficient-domain residual
\[
  (B_3-\mathbf 1_n r_1)\Had\bigl(\mA\vecu-B_0-B_1(\vecm\Vert\veck)-B_2r_0\bigr)-1.
\]
Injectivity then forces this coefficient-domain residual itself to vanish, which is exactly the displayed
identity.
\end{proof}

\begin{corollary}[Abstract BB-tran witness consequence]\label{cor:showing-to-abstract}
Under the assumptions of Theorem~\ref{thm:replay-to-coeff}, the extracted witness additionally satisfies
\[
  (B_3-\mathbf 1_n r_1)\Had Z=1,
  \qquad
  \mA\vecu=h_{\mathrm{tran},\vecm\Vert\veck,(r_0,r_1)}(B),
  \qquad
  \prftag=\PRF(\veck,\nonce).
\]
\end{corollary}

\begin{proof}
The inverse equation $(B_3-\mathbf 1_n r_1)\Had Z=1$ is part of the compiled relation.  By
Theorem~\ref{thm:replay-to-coeff},
\[
  (B_3-\mathbf 1_n r_1)\Had\bigl(\mA\vecu-B_0-B_1(\vecm\Vert\veck)-B_2r_0\bigr)=1.
\]
Combining this with the inverse equation identifies the bracketed term with $Z$, hence
\[
  \mA\vecu=B_0+B_1(\vecm\Vert\veck)+B_2r_0+Z
  =h_{\mathrm{tran},\vecm\Vert\veck,(r_0,r_1)}(B).
\]
The PRF equation is already part of the showing relation.
\end{proof}

\subsection{PRF companion relation}
\label{subsec:prf-checkpointed}
The public tag is tied to the same hidden key carried in the signed message through a companion PRF
relation.  Selector constraints bind the PRF key lanes to the key rows $\veck$ and enforce
$\veck\in\mathcal K_{\mathsf{key}}$; checkpoint constraints certify the nonlinear Poseidon2-style S-box steps;
linear layers and round constants are replayed from public parameters; and the final feed-forward and truncation
constraints bind the public pair $(\nonce,\prftag)$.

A crucial invariant is that the PRF key rows are not separate from the signature/message rows: the compiled
showing statement is a \emph{single} committed witness family.  If auxiliary alias rows are introduced for
engineering convenience, they must be tied back to the original key rows by explicit copy constraints.  This
prevents a prover from using one hidden key for the BB-tran signature relation and another hidden key for the
PRF relation.

\subsection{Statement strength}
\label{subsec:statement-strength}
The SmallWood proof hides a family of row polynomials, not merely a tuple of scalars.  Random masks on
$\Omega'$ statistically hide unopened row values, while the DECS/LVCS/PCS stack reveals only the evaluations
needed to check the PACS constraints.  Consequently, the verifier learns only the public statement: during
issuance this is $(\vecc,r_0^I,r_1^I,T)$ together with public parameters, and during showing it is
$(\nonce,\prftag)$ together with the public verification context.

The replay-space statement of Theorem~\ref{thm:replay-to-coeff} is therefore the only statement certified
directly by the post-sign proof.  The coefficient-domain BB-tran identity used in the security arguments is not
assumed separately; it is a theorem-level consequence of the compiled proof relation and the exact bridge
constraints.

\begin{corollary}[Consequence of an accepting post-sign proof]\label{cor:showing-proof-consequence}
If the post-sign NIZK is knowledge sound for the compiled showing relation, then every accepting proof yields a
witness
\[
  (\vecu,\vecm,\veck,r_0,r_1,Z)
\]
satisfying the coefficient-domain BB-tran signature relation and the PRF-tag relation.  In particular,
\[
  \mA\vecu=h_{\mathrm{tran},\vecm\Vert\veck,(r_0,r_1)}(B)
  \qquad\text{and}\qquad
  \prftag=\PRF(\veck,\nonce).
\]
\end{corollary}
